

\section*{Chapitre 1 - Nombres entiers}
\addcontentsline{toc}{section}{Chapitre 1 - Nombres entiers}


\subsection*{I. Écrire les nombres entiers}
\addcontentsline{toc}{subsection}{I. Écrire les nombres entiers}

\begin{Définition*}
Pour écrire les nombres, on utilise 10 symboles appelés chiffres : 0, 1, 2, 3, 4, 5, 6, 7, 8, 9.
\end{Définition*}


\begin{Remarque}
Pour faciliter la lecture des grands nombres, on regroupe les chiffres par trois en partant de la droite : 
le nombre 3562382598 s'écrit 3 562 382 598.
\end{Remarque}


\textbf{Règles d'orthographe pour l'écriture des nombres en lettres}

\begin{itemize}
\item Si un nombre s'écrit en plusieurs mots, on met un trait d'union entre chaque mot.
\item Les mots « vingt » et « cent » se terminent par un « s » s'ils sont le dernier mot du nombre et qu'ils sont multipliés (six-cents, quatre-vingts, mais six-cent-dix et quatre-vingt-cinq).
\item Le mot « mille » est invariable, il ne prend jamais de « s ».
\item Les mots « million » et « milliard » s'accordent au pluriel.
\end{itemize}



\begin{Exemple}
3 562 382 598 s'écrit \ligne

\ligne

\vspace{-0.2cm}
\end{Exemple}

\vspace{-0.2cm}
\subsection*{II. Ranger les nombres entiers}
\addcontentsline{toc}{subsection}{II. Ranger les nombres entiers}



%\begin{Définition}
%Un \textbf{milliard} est égal à mille millions. \\
%On peut alors écrire 1 000 000 000 = 1 000 000 $\times$ 1 000.
%\end{Définition}

\begin{Propriété}[c]
$ $ \vspace{-0.1cm}
\begin{center}
\arrayrulecolor{Couleur1}
\scalebox{0.8}{%
\begin{tabular}{|*{12}{>{\centering\arraybackslash}p{0.9cm}|}}
\hline
\multicolumn{3}{|c|}{Classe des milliards} & \multicolumn{3}{c|}{Classe des millions} & \multicolumn{3}{c|}{Classe des milliers} & \multicolumn{3}{c|}{Classe des unités} \\
\hline
\rotatebox{90}{\shortstack{Centaines\\de milliards}} & \rotatebox{90}{\shortstack{Dizaines\\de milliards}} & \rotatebox{90}{\shortstack{Unités\\de milliards}} & 
\rotatebox{90}{\shortstack{Centaines\\de millions}} & \rotatebox{90}{\shortstack{Dizaines\\de millions}} & \rotatebox{90}{\shortstack{Unités\\de millions}} & 
\rotatebox{90}{\shortstack{Centaines\\de milliers}} & \rotatebox{90}{\shortstack{Dizaines\\de milliers}} & \rotatebox{90}{\shortstack{Unités\\de milliers}} & 
\rotatebox{90}{Centaines} & \rotatebox{90}{Dizaines} & \rotatebox{90}{Unités} \\
\hline
 &  &  &  &  &  &  &  &  &  &  &  \\
\hline
\end{tabular}
}
\end{center}

\end{Propriété}


\begin{Exemple}

En t'aidant du tableau ci-dessus, réponds aux questions suivantes : 
\begin{enumerate}
\item Dans le nombre 3 562 382 598, quel est le rang du chiffre 6 ?

\ligne

\item Dans le nombre 3 562 382 598, quel est le nombre des dizaines de milliers ? 

\ligne

\item Dans  3 562 382 598, déterminer le nombre de dizaines. 

\ligne

\item Décomposer 3 562 382 598 selon ses rangs. 

\ligne

\ligne

\item En déduire son chiffre des centaines de milliers.

\ligne

\end{enumerate}
\vspace{-0.2cm}
\end{Exemple}


